\documentclass[11pt,a4paper]{article}
\usepackage[T1]{fontenc}
\usepackage[utf8]{inputenc}
\usepackage[margin=1in]{geometry}
\usepackage{amsmath,amssymb,amsthm}
\usepackage{stmaryrd}
\usepackage{booktabs}
\usepackage{tabularx}
\usepackage{array}
\usepackage{enumitem}
\usepackage{listings}
\usepackage{xcolor}
\usepackage[expansion=false]{microtype}
\usepackage{tikz}
\usetikzlibrary{arrows.meta,positioning,fit,calc}
\usepackage[htt]{hyphenat}
\usepackage[hidelinks,pdftitle={F1R3Comb-Mat v0.2}]{hyperref}
\setlength{\emergencystretch}{3em}

\newcommand{\code}[1]{\texttt{#1}}
\newcommand{\Comb}{\textsc{F1R3Comb}}
\newcommand{\Mat}{\textsc{F1R3Comb-Mat}}
\newcommand{\Kind}{\textsc{K1ndl1ng}}
\newcommand{\Camp}{\textsc{CampF1R3}}
\newcommand{\Bell}{\textsc{B3ll0ws}}
\newcommand{\Fx}{\textsc{F1R3x}}
\newcommand{\Rhobots}{\textsc{Rhobots}}
\newcommand{\MUST}{\textsc{must}}
\newcommand{\MUSTNOT}{\textsc{must not}}
\newcommand{\SHOULD}{\textsc{should}}
\newcommand{\SHOULDNOT}{\textsc{should not}}
\newcommand{\MAY}{\textsc{may}}

\newcommand{\quo}[1]{@#1}
\newcommand{\drp}[1]{{*}#1}
\newcommand{\nil}{\mathbf{0}}
\newcommand{\nequiv}{\equiv_{\mathcal{N}}}
\newcommand{\nms}[1]{\mathcal{N}^{*}(#1)}
\newcommand{\Nsort}{\mathsf{N}}
\newcommand{\Psort}{\mathsf{P}}
\newcommand{\red}{\longrightarrow}
\newcommand{\wred}{\longrightarrow^{*}}
\newcommand{\scong}{\equiv}
\newcommand{\comps}[1]{\lvert #1 \rvert}

\newcommand{\mm}{\mathsf{m}}
\newcommand{\dd}{\mathsf{d}}
\newcommand{\kk}{\mathsf{k}}
\newcommand{\fw}{\mathsf{fw}}
\newcommand{\bl}{\mathsf{b}_{\mathsf{l}}}
\newcommand{\br}{\mathsf{b}_{\mathsf{r}}}
\newcommand{\sy}{\mathsf{s}}
\newcommand{\ev}{\mathsf{e}}
\newcommand{\qq}{\mathsf{q}}
\newcommand{\cons}{\mathsf{cons}}
\newcommand{\consm}{\mathsf{cons}_{\mm}}
\newcommand{\consd}{\mathsf{cons}_{\dd}}
\newcommand{\conss}{\mathsf{cons}_{\sy}}
\newcommand{\conse}{\mathsf{cons}_{\ev}}
\newcommand{\consf}{\mathsf{cons}_{\fw}}

\newcommand{\Spec}{\mathcal{S}}
\newcommand{\Nm}{\mathcal{N}}
\newcommand{\cre}[1]{a^{\dagger}_{#1}}
\newcommand{\ann}[1]{a_{#1}}
\newcommand{\Wt}{\mathsf{W}}

\theoremstyle{definition}
\newtheorem{definition}{Definition}[section]
\newtheorem{requirement}[definition]{Requirement}
\newtheorem{decision}[definition]{Decision}
\newtheorem{obligation}[definition]{Obligation}
\newtheorem{hazard}[definition]{Hazard}
\newtheorem{erratum}[definition]{Erratum}
\newtheorem{remark}[definition]{Remark}
\newtheorem{example}[definition]{Example}
\theoremstyle{plain}
\newtheorem{proposition}[definition]{Proposition}
\newtheorem{lemma}[definition]{Lemma}

\newcommand{\Lft}{\mathsf{L}}
\newcommand{\Rgt}{\mathsf{R}}
\newcommand{\inst}{\mathsf{inst}}
\newcommand{\ERECT}{\mathrm{ERECT}}
\newcommand{\ALLOC}{\mathrm{ALLOC}}
\newcommand{\enc}[1]{\llbracket #1 \rrbracket}
\newcommand{\inp}[3]{\mathsf{for}(\,#1 \leftarrow #2\,)#3}
\newcommand{\out}[2]{#1!(#2)}
\newcommand{\Fam}{\mathcal{F}}

\lstset{
  basicstyle=\ttfamily\footnotesize,
  columns=fullflexible,
  keepspaces=true,
  breaklines=true,
  frame=leftline,
  framesep=6pt,
  xleftmargin=8pt,
  rulecolor=\color{black!30},
  showstringspaces=false,
}

\title{\textbf{\Mat}\\[1mm]
\large The rho combinators as sparse linear algebra:\\
a Rust implementation of the matrix representation,\\
the bulk-synchronous machine, and the transfer operator\\[2mm]
\normalsize Specification v0.2 --- companion to \Comb{} v0.4}
\author{F1R3FLY.io}
\date{23 September 2026}

\begin{document}
\maketitle

\begin{abstract}
\noindent
This is the second version of the specification for \Mat, the arrow from a
\Comb{} term to the matrix representation of \cite{matnote} and the
bulk-synchronous machine that runs it: the route from the rho calculus,
through the combinators, to a data-parallel device. Since v0.1, draft~3 has
changed the translation it runs \cite[\S6]{draft3}. Apparatus names are no
longer positional, because positional names are unsound for code instantiated
twice; a first phase compiles into the combinators with a binder, and a second
eliminates the binder by allocating each run-time instance's apparatus beneath
an \emph{address} it receives, and erecting its atoms there by a constructor
followed by $\ev$. The four constructors have become a family, one per atom
shape. \Comb{} has followed, to v0.4 \cite{combspec}.

For a machine whose whole purpose is to run many things at once, the
decisive part of this is phase two, and v0.2 is organised around it. Three
findings, each backed by the reference implementation, shape the
specification. First, concurrent instances erected by the literal
construction \cite[Def.~6.8]{draft3} mix their apparatus in most runs, and
maximal progress --- the device's native schedule --- mixes more than
sequential choice does (\cite[Find.~14.1]{matnote}: 58\% of two-instance runs
against 32\%, every run from eight instances). Second, the discipline draft~3
proves sufficient for correctness \cite[Def.~6.10, Lem.~6.11]{draft3} cannot
be met with fixed-arity constructors for any atom mentioning two allocated
names --- which includes the distributor and the gate --- so the conforming
construction \Comb{} v0.4 mandates \cite[Req.~5.16]{combspec} does not exist
\cite[Prop.~14.3]{matnote}. Third, a single-input, multi-hole constructor,
which is draft~3's own context-indexed alternative \cite[Rem.~3.14]{draft3}
promoted to a primitive, meets the discipline by construction, erects an
instance in two steps, keeps rows four wide, and runs the recursion
combinator at a constant 7 parallel steps and 16 new names per unfolding
however deep its address (\cite[Find.~14.4]{matnote}).

v0.2 therefore supports both erections: the literal one, because it is what
\Comb{} emits today and the machine must be able to show what it does, and
single-input instantiation, behind a feature, as the construction the device
path is built for. It generates its rule table from \Comb's shape table, makes
the row width a per-artefact parameter, adopts \Comb's address allocator and
freshness checks, adds an instance-provenance column for tracing and for a
cross-instance resolver, and records that the operator layer's closure applies
to no compiled program with an input, since every one contains $\ev$. The
three errata v0.1 raised against the notes are discharged by their revisions;
the questions this version puts to the papers are about phase two.
Key words \MUST, \MUSTNOT, \SHOULD, \MAY{} are used in the sense of RFC~2119.
\end{abstract}

\tableofcontents

%======================================================================
\section{Purpose and scope}\label{sec:purpose}

\subsection{The route, and where this document sits on it}

The programme is a compiler from the rho calculus to a data-parallel device
by way of the combinators and their matrix reading. It has four stages, and
this document specifies the third.
\begin{center}
\begin{tabularx}{\linewidth}{@{}l>{\raggedright\arraybackslash}Xl@{}}
\toprule
Stage & What it does & Specified in\\
\midrule
1 & K0 text to a canonical normal form & \cite{k1ndl1ng}\\
2 & Normal form to combinator term, in two phases & \Comb{} v0.4 \cite{combspec}\\
3 & Combinator term to matrix representation; the bulk-synchronous machine;
    the transfer operator & this document\\
4 & The machine's kernels on a device & Section~\ref{sec:device}, as a
    contract\\
\bottomrule
\end{tabularx}
\end{center}

\noindent
Two properties are fixed from the start and govern everything below. The
machine must compute exactly one step per state, seed and configuration, so
that a host and a device can be compared trace for trace
(\cite[P4]{matnote}); and it must run what stage two emits, including every
instance of every unit, at the width the program has
(\cite[Princ.~10.4]{matnote}).

\subsection{What exists today}

\begin{center}
\begin{tabularx}{\linewidth}{@{}l>{\raggedright\arraybackslash}X@{}}
\toprule
Artefact & What it gives this project\\
\midrule
\code{rho-combinators-draft3} & The calculus and its two presentations
(\S3.2--3.3); the constructor family (Def.~3.8, Rem.~3.9); erection at a
computed name (Rem.~3.11); the context-indexed constructor (Rem.~3.14); the
positional-naming bug (\S6.1); phase one (\S6.2) and phase two (\S6.3,
Def.~6.6, Def.~6.8, Ex.~6.9); the discipline and the permutation lemma
(Def.~6.10, Lem.~6.11, Rem.~6.12); freshness by size (Lem.~6.14); the weighted
atom count (Lem.~9.2, Prop.~9.9)\\
\code{f1r3comb-spec-v04} & The term type and shape table; tags; the
intermediate representation; phase two as a compiler; the address allocator;
the golden terms $W$ and $D_{x}$; the diagnostics this document shares\\
\code{rhocomb-matrix} v1.2 & The state space and incidence (\S4, \S7), the
step (\S8), the operator (\S5--6), the plan (\S11), and \S14: phase two seen
from the machine, the mixing measurements, the multiplicity proposition,
context instantiation\\
\code{rhocombmat.py} v1.2, \code{experiments.py}, \code{phase2.py} & The
reference machine with a generated rule table and the context constructor;
E1--E4 (\code{results.txt}); the phase-two experiments
(\code{phase2-results.txt})\\
\code{rhocomb-logic} rev.~2 & Components as a complete invariant (Def.~2.9);
drop elimination (Lem.~2.8); injective quotation (Prop.~2.10); the weighted
count as a formula (Def.~4.2, Prop.~6.17)\\
\code{campf1r3} (the repository) & The workspace; \code{Hash32} and
\code{blake2b\_256}; no external dependencies\\
\bottomrule
\end{tabularx}
\end{center}

\subsection{Nomenclature}\label{sec:nomen}

\begin{center}
\begin{tabularx}{\linewidth}{@{}l>{\raggedright\arraybackslash}X@{}}
\toprule
Term & Meaning in this document\\
\midrule
\emph{term} & A \code{f1r3comb\_term::Term}: a closed drop-free normal form
under presentation~A, a flat multiset of atoms \cite[Req.~4.3]{combspec}\\
\emph{shape table} & \Comb's table of atom shapes with arities and argument
sorts, from which constructors and rules are generated
\cite[Req.~4.5]{combspec}\\
\emph{family}, $\Fam(P)$ & The constructors a program uses
\cite[Req.~4.7]{combspec}, recorded in the artefact header\\
\emph{row width}, $w$ & The largest arity in the program's signature; the
fixed width of every row in one run\\
\emph{name id} & A \code{NameId(u32)}: the in-run index of a name modulo
$\nequiv$\\
\emph{row}, \emph{row ref} & One atom occurrence; its position
\code{(Shape, u32)} in its shape table\\
\emph{address}, \emph{leaf} & The name an instance receives and allocates
beneath; a name $\sigma\!\cdot\!w$ so allocated \cite[Def.~6.6]{draft3}\\
\emph{static channel} & A name in the phase-two image, shared by every
instance \cite[\S1.3]{combspec}\\
\emph{unit}, \emph{instance} & One quotation's phase-two image; one
run-time copy of it at one address\\
\emph{provenance} & The instance a row belongs to, as far as the machine can
tell; machine metadata, never semantics (Section~\ref{sec:prov})\\
\emph{erection} & Building an atom at computed names: a constructor, then
$\ev$ \cite[Rem.~3.11]{draft3}; \emph{literal} per \cite[Def.~6.8]{draft3},
or by \emph{instantiation} per Section~\ref{sec:inst}\\
\emph{redex}, \emph{step}, \emph{width} & As in v0.1: a rule with the row refs
it consumes; one maximal independent set; its size\\
\bottomrule
\end{tabularx}
\end{center}

\subsection{Goals}

\begin{enumerate}[leftmargin=1.6em]
\item A total, lossless, round-trippable map from \Comb{} terms to the matrix
representation, and back, for every shape \Comb's shape table can generate.
\item Redex finding as joins derived from the generated rule table, with three
independent finders that agree.
\item A bulk-synchronous step defined bit-exactly, so host and device traces
are comparable.
\item Phase-two images run correctly: \Comb's address allocator and freshness
checks, per-instance provenance, and both erections --- literal, to test what
\Comb{} emits, and single-input instantiation, to build on.
\item The instruments that gate device work: width, with erection separated
from computation; address growth; cross-instance mixing.
\item On the destructor-free sublattice, the transfer operator, its closure
and the pre-image primitive, scoped honestly (Section~\ref{sec:op}).
\item \Comb's portability envelope: no external dependencies,
\code{forbid(unsafe\_code)}, WebAssembly-clean, bit-identical across hosts
and core counts.
\end{enumerate}

\subsection{Non-goals}\label{sec:nongoals}

Device code (Section~\ref{sec:device} is a contract, and gated); presentation
B (Decision~\ref{dec:aonly}); weights; the logic; the history unfolding;
intern-table collection; phase one and phase two themselves, which are
\Comb's. \Mat{} takes a \code{Term} and \MUSTNOT{} depend on
\code{f1r3comb-lower} or \code{f1r3comb-ir}.

%======================================================================
\section{Position relative to \Comb}\label{sec:position}

\subsection{Two machines, one term type}

\begin{center}
\begin{tabularx}{\linewidth}{@{}l>{\raggedright\arraybackslash}X>{\raggedright\arraybackslash}X@{}}
\toprule
& \Comb{} machine (\code{f1r3comb-run}) & \Mat{} (\code{f1r3comb-par})\\
\midrule
Input & \code{Term} and artefact header & \code{Term} and artefact header\\
State & \Camp{} store keyed by \code{Hash32} & shape tables of \code{NameId}
rows, width $w$\\
Name identity & content hash of the encoding & structural intern in-run,
content hash at the boundary\\
Event & one comm & one step, a maximal independent set\\
Resolver & store order key; an adversarial profile for testing
\cite[Obl.~11.2]{combspec} & priority function; a cross-instance profile
(Section~\ref{sec:resolvers})\\
Rules & generated from the shape table & generated from the same table\\
Addresses & allocator \cite[Req.~9.5]{combspec} & the same allocator\\
Output & a scene & a scene\\
\bottomrule
\end{tabularx}
\end{center}

\subsection{Decisions}

\begin{decision}[Presentation A only]\label{dec:aonly}
\Mat{} implements presentation~A, and under \Comb's \code{presentation-b}
feature its crates compile to stubs returning \code{comb-mat-presentation}.
\Comb{} v0.4 withdrew one of its own reasons for A \cite[Rem.~2.2]{combspec};
this decision never rested on it. It rests on row width: under B a message
payload is a process, rows are ragged, and (G1) of \cite[\S2.1]{matnote}
fails. The A/B comparison \Comb{} requires \cite[Obl.~11.8]{combspec} is run
on \code{f1r3comb-run}.
\end{decision}

\begin{decision}[A store row holds the quotation of its payload]
\label{dec:qrow}
Unchanged from v0.1, and now with a second dividend. The term type keeps
$\qq(a,p)$ with a process payload \cite[Dec.~2.1]{combspec}; the row is
$(\qq,\mathrm{id}(a),\mathrm{id}(\quo{p}))$, lossless because $\quo{-}$ is
injective on normal forms \cite[Prop.~2.10]{logic}, and the shape tag carries
the $\mm$/$\qq$ distinction the logic needs. In the row representation the
member $\cons_{\qq}$ \cite[Def.~3.8]{draft3} is uniform with every other: its
conclusion is the row $(\qq,v_{1},v_{2})$, and A's coercion
$\quo{\qq(v_{1},\drp{v_{2}})}$ \cite[Rem.~3.9]{draft3} never appears
(\cite[Rem.~3.1]{matnote}). \cite[Req.~4.8]{combspec}'s requirement that the
two readings agree holds by construction here and \MUST{} be tested on
\code{f1r3comb-run}'s side only.
\end{decision}

\begin{decision}[Decode by arena read]\label{dec:decode}
Unchanged from v0.1. \cite[Req.~9.3]{combspec} forbids the \Camp{} machine to
project a payload out of a name, because the short-cut is unsound under B.
\Mat{} is A-only and interns drop-free components, so its $\mathsf{eval}$ is
an arena read, tested against the smart-constructed $\drp{v}$.
\end{decision}

\begin{decision}[The rule table is generated from \Comb's shape table]
\label{dec:ruletable}
v0.1 asked for one rule table shared by both machines. \Comb{} v0.4 goes
further and generates it \cite[Req.~4.5]{combspec}. \Mat{} \MUST{} consume the
same generator, living in \code{f1r3comb-term}, and \MUSTNOT{} list rules of
its own. Join plans, rule classes, the grading and the row width are derived
from the generated table (Section~\ref{sec:rules}).
\end{decision}

\begin{decision}[Two erections, one of them behind a feature]
\label{dec:twoerect}
\Mat{} \MUST{} run terms erected literally per \cite[Def.~6.8]{draft3}, since
that is what \Comb{} emits, and \MUST{} support the context constructor
$\inst$ of Section~\ref{sec:inst} behind the feature \code{inst}, off by
default until the paper settles phase two's primitive
(Obligation~\ref{obl:inst}). The feature adds one shape and one rule; the
rest of the machine is indifferent to it. The device path is specified for
the \code{inst} construction (Section~\ref{sec:device}), because it is the one
that is both correct under concurrency and fixed-width.
\end{decision}

%======================================================================
\section{Phase two, as the machine meets it}\label{sec:phase2}

This section fixes what \Mat{} must do with a phase-two image, and why. The
argument is in \cite[\S14]{matnote}; what follows is its consequences as
requirements.

\subsection{Concurrency is the normal case}

Two deploys of one program share every static channel, and so do two
unfoldings of one recursive unit. Draft~3's correctness rests on the
permutation lemma \cite[Lem.~6.11]{draft3}, which holds for images meeting the
single-address-input discipline. The literal erection does not meet it
\cite[Req.~5.15]{combspec}, and \cite[Find.~14.1]{matnote} measures the
consequence on the reference machine:

\begin{center}\small
\begin{tabular}{@{}lrrrr@{}}
\toprule
instances & 2 & 4 & 8 & 16\\
\midrule
maximal progress: runs with a mixed atom (of 50) & 29 & 48 & 50 & 50\\
sequential uniform: runs with a mixed atom (of 50) & 16 & 38 & 49 & 50\\
single-input instantiation, either scheduler & 0 & 0 & 0 & 0\\
\bottomrule
\end{tabular}
\end{center}

\noindent
A bulk-synchronous machine fires every enabled constructor at once, so it
exhibits the cross whenever two instances' leaves are present together, which
on a device is always. This is why \Mat{} treats mixing as something to
measure, not something the schedule might avoid.

\subsection{The conforming construction does not exist}

\Comb{} v0.4 responds by mandating a two-stage erection: a single-input
bootstrap that erects, at address-derived channels, erectors that may then
join freely \cite[Req.~5.16]{combspec}. \cite[Prop.~14.3]{matnote} shows that
under the discipline in the form the permutation lemma uses --- no redex fired
at a static channel draws occurrences of one instance's address from two
premises --- no reachable atom contains the address twice. An erector that
joins two leaves contains the address twice, and so does every atom
mentioning two allocated names, which includes the distributor
$\dd(q_{1},p_{1},q_{2})$ and the gate $\sy(p_{0},b,c)$. Stage~(ii) of the
two-stage scheme is therefore unreachable from stage~(i), and draft~3's
\cite[Rem.~6.12]{draft3} describes no construction.

The obstruction is arity, not routing. A fixed-arity constructor places each
input once, so building a name that mentions the address twice from one copy
of it requires duplicating and rejoining, and the rejoin is the forbidden
join. A constructor that places one input at many holes does not.

\begin{hazard}[Correctness of compiled programs is gated]\label{haz:gate}
Until phase two has a primitive that meets the discipline, a compiled program
that runs two instances of a unit concurrently has no correctness guarantee on
any machine, and on this one fails in most runs. Every width measurement on
compiled programs (Section~\ref{sec:instr}) \MUST{} be labelled with the
erection used, and results obtained with literal erection \MUST{} be treated
as measurements of the schedule, not of the program.
\end{hazard}

\subsection{Single-input context instantiation}\label{sec:inst}

\begin{definition}[The context constructor]\label{def:inst}
For a context $C$ --- a term with a hole in any number of name positions,
including beneath quotation ---
\[
  \inst(a,f,C) \mid \mm(a,v) \;\red\; \mm(f,\quo{C[v]}).
\]
A unit's phase-two image under \code{inst} is $\inst(r,f,C_{\mathrm{unit}})
\mid \ev(f)$, where $C_{\mathrm{unit}}$ is the unit's phase-one image with
each bound name replaced by the leaf at a fixed-length path beneath the hole,
and with the spare leaves posted for nested units and self-instantiation
\cite[Req.~5.11]{combspec}.
\end{definition}

This is \cite[Rem.~3.14]{draft3}'s context-indexed constructor, with one
premise. It meets the discipline by construction (one address input at the
static $r$); permuting addresses among $\inst$ atoms is the harmless
permutation of \cite[Lem.~6.11]{draft3}; and it removes the allocation tree,
the per-node gadget channels \cite[Req.~5.9]{combspec}, the fan-out count
\cite[Req.~5.12]{combspec} and the constructors of constructors. Freshness is
unchanged: leaves exceed the address in size, and the address and statics
exceed every name of the image \cite[Lem.~6.14]{draft3},
\cite[Req.~5.22]{combspec}.

\begin{requirement}[The context is a name with a reserved hole]\label{req:hole}
$C$ \MUST{} be represented as an interned name in which a reserved
\code{NameId::HOLE} stands for the hole. \code{HOLE} \MUST{} have no
counterpart in the term type other than inside the third argument of an
$\inst$ atom, \MUST{} be rejected by import anywhere else
(\code{comb-mat-hole}), and \MUST{} be unconstructible: no rule's conclusion
may contain it. In the term encoding it \MUST{} take a tag in the reserved
block, allocated by \Comb{} (Section~\ref{sec:paper}, item~4).
\end{requirement}

\begin{requirement}[Instantiation is a memoised substitution]\label{req:subst}
The encode of $\inst$ \MUST{} be $C\{v/\code{HOLE}\}$ computed over the
interned form, hereditarily beneath quotation, with a per-firing memo keyed
by name id, so that its cost is the number of distinct names in $C$, not the
size of $C$'s tree. Its result names \MUST{} be interned in canonical order
(Requirement~\ref{req:insertorder}).
\end{requirement}

\begin{requirement}[The measured cost is a regression]\label{req:dxlaw}
\cite[Find.~14.4]{matnote} runs $D_{x}$ of \cite[Ex.~6.9]{draft3} under
\code{inst}: exactly 7 parallel steps and 16 new names per unfolding for 200
unfoldings, and between 9 and 11 sequential steps (mean 10, the draft's 2
erection steps plus 8). \code{f1r3comb-check} \MUST{} reproduce these as laws
over the number of unfoldings, to $10^{4}$ unfoldings
(Requirement~\ref{req:targets}).
\end{requirement}

\subsection{Provenance}\label{sec:prov}

A trace of a phase-two run is unreadable without knowing which instance each
atom belongs to \cite[Req.~9.5]{combspec}, and a cross-instance resolver needs
the same information. It is available to the machine without inspecting
names.

\begin{definition}[Provenance]\label{def:prov}
Every row carries a provenance tag \code{Prov(u32)}: \code{STATIC} for rows
imported from the artefact, and for the rest the tag of the root address
beneath which the row's allocated names lie, computed when the row is
produced as follows. A row produced by a redex takes the tag of the redex's
premises if they agree on one non-\code{STATIC} tag, \code{STATIC} if all
are \code{STATIC}, and \code{MIXED} otherwise; a row whose names include a
leaf takes that leaf's root's tag, found in $O(1)$ from the intern table's
root column (Requirement~\ref{req:roots}).
\end{definition}

\begin{requirement}[Provenance is metadata]\label{req:provmeta}
Provenance \MUSTNOT{} influence which redexes exist, which are selected under
\code{MaximalProgress}, or what they produce. It \MAY{} influence selection
only under the resolver of Section~\ref{sec:resolvers} built for it. It
\MUST{} appear in traces, rendered as the root address's content hash, and a
row tagged \code{MIXED} \MUST{} raise a trace event at level \code{races} or
above. Under \code{inst} no row \SHOULD{} ever be tagged \code{MIXED}, and
the test suite \MUST{} assert it.
\end{requirement}

%======================================================================
\section{The representation}\label{sec:rep}

\subsection{Names, the intern table, and addresses}

\begin{lstlisting}
#[derive(Copy, Clone, PartialEq, Eq, PartialOrd, Ord, Hash, Debug)]
pub struct NameId(pub u32);            // NameId::HOLE reserved (Requirement 3.3)

pub struct Row<const W: usize> { pub shape: ShapeId, pub args: [NameId; W], pub prov: Prov }

pub struct InternTable {
    arena:  Vec<(ShapeId, SmallArgs)>,  // components of every name, contiguous
    span:   Vec<(u32, u32)>,            // NameId -> (offset, len)
    key:    HashMap<Box<[..]>, NameId, FixedHasher>,
    size:   Vec<u64>,                   // encoding size, for freshness by size
    root:   Vec<Option<NameId>>,        // the address a leaf lies beneath
    depth:  Vec<u32>,                   // quotation depth, for canonical export
    chash:  Vec<Option<Hash32>>,        // content hash, lazily
    budget: NameBudget,
}
\end{lstlisting}

\begin{requirement}[Structural interning; the hot path is hash-free]
\label{req:hashfree}
As v0.1: a name's key is its sorted component rows over child ids, \code{quote}
and \code{drop} compute no BLAKE2b, and the content hash of
\cite[Req.~4.10]{combspec} is computed only at the boundary. The arena stores
rows of the name's own width, so names built from shapes of different arity
share one table.
\end{requirement}

\begin{remark}[Structural interning is the prefix sharing \Comb{} asks for]
\cite[Req.~9.7]{combspec} asks that descending one address level cost $O(1)$
rather than $O(|\sigma|)$, and notes that without it a recursive program
degrades quadratically \cite[Req.~13.1]{combspec}. Structural interning
provides it by construction: $\Lft\sigma$ is one new entry whose single
component refers to $\sigma$'s id, and \cite[Find.~14.4]{matnote} measures 16
new names per unfolding of $D_{x}$ at every depth. The content hash, being
of the full encoding, is $O(\text{depth})$ per new address; this is the
reason, beyond v0.1's, that \Mat{} keeps it off the hot path, and the
reason for item~7 of Section~\ref{sec:paper}.
\end{remark}

\begin{requirement}[Size, root and depth columns]\label{req:roots}
\code{size} \MUST{} hold each name's encoding length, computed from its
children in $O(\text{components})$ at interning time, so that the freshness
test by size \cite[Lem.~6.14]{draft3} is an integer comparison. \code{root}
\MUST{} hold, for a leaf, the root address beneath which it lies, computed at
interning time from the child's entry when the name has the $\Lft$ or $\Rgt$
shape of \cite[Def.~6.6]{draft3}. Both columns are machine metadata subject
to Requirement~\ref{req:noidleak}.
\end{requirement}

\begin{requirement}[Ids never leave the process]\label{req:noidleak}
As v0.1: no trace, report, diagnostic or export \MAY{} contain a \code{NameId}
or row position, except the canonically relabelled \code{.cmat}. Externally a
name is its \code{Hash32}.
\end{requirement}

\begin{requirement}[Budget]\label{req:budget}
\cite[Req.~9.7]{combspec}'s budget: maximum distinct names, maximum
components per name, maximum arena length, and maximum address depth.
Exceeding any \MUST{} halt the step before it commits, with
\code{comb-name-budget} naming the rule and, for an address, the unit.
Because every loop descends \cite[Haz.~9.6]{combspec}, the budget is an
operational parameter and \MUST{} be reported with its high-water marks in
every run report.
\end{requirement}

\subsection{State}

\begin{requirement}[Structure of arrays, width $w$]\label{req:soa}
A state is one table per shape in the program's signature, each holding
$\mathrm{ar}(A)$ columns of \code{NameId} plus a provenance column. Column
arrays are contiguous; the incidence $M^{A}_{j}$ is the column array
(\cite[Def.~7.1]{matnote}). The width $w$ \MUST{} be read from the artefact
header, \MUST{} be a const-generic parameter of the machine instantiated for
$w \in \{4,5,6\}$, and an artefact whose signature needs more \MUST{} be
rejected with \code{comb-mat-width}. Under \code{inst} $w = 4$ for every
program: the base shapes have arity at most three, $\cons_{\dd}$ and
$\cons_{\sy}$ four, and $\inst$ three.
\end{requirement}

\begin{remark}[Why the width is not a constant any more]
Under literal erection phase two erects every phase-one atom that mentions a
bound name, and the (o5) chains of phase one are constructor atoms on bound
names, so a program with an (o5) occurrence needs $\cons_{\cons_{\dd}}$, of
arity five with four premises (\cite[\S14.1]{matnote}). \Comb{} v0.4's
family computation includes these \cite[Req.~4.7]{combspec}, but its tag
table does not (Section~\ref{sec:paper}, item~3). The two-stage erection, were
it constructible, would add a third level.
\end{remark}

\subsection{Import, export, and the canonical format}

\begin{requirement}[Import and export]\label{req:import}
As v0.1, extended to the generated signature: import flattens components in
encoding order and interns bottom-up with an explicit stack; export rebuilds
through \code{f1r3comb-term}'s smart constructors; import--export is the
identity on encodings and export--import on markings. Provenance is
\code{STATIC} on import and is not exported.
\end{requirement}

\begin{requirement}[\code{.cmat}]\label{req:cmat}
As v0.1, with the header extended by the fields \Comb's artefact header
carries \cite[\S8.1]{combspec}: presentation, feature set (including
\code{inst}), erection construction, $\Fam(P)$, $w$, the root-address size
bound, and the static channels. Canonical relabelling orders names by
quotation depth then content hash; \code{HOLE} is relabelled to the reserved
index $0$. For terms $P \scong Q$ the encodings \MUST{} be byte-equal.
\end{requirement}

%======================================================================
\section{The rule table}\label{sec:rules}

\subsection{Generation}

\begin{lstlisting}
pub enum Premise { Loud { slot: u8 }, Quiet { slot: u8 } }
pub enum Conclusion {
    Atoms(Vec<(ShapeId, Vec<Var>)>),     // class (a)
    Encode { deliver: Var, node: Node }, // class (b): m(deliver, @node)
    Decode { payload: Var },             // class (c): *payload
}
pub enum Node { Union(Var, Var), Atom(ShapeId, Vec<Var>), Instantiate { ctx: Var, with: Var } }
pub fn rules(shapes: &ShapeTable, features: Features) -> RuleTable;   // in f1r3comb-term
\end{lstlisting}

\begin{requirement}[What the generator emits]\label{req:gen}
From the shape table the generator \MUST{} emit: the seven routing and
destructor rules and $\mathsf{quote}$ of \cite[Def.~3.3]{draft3};
$\mathsf{make}_{\mid}$; for each constructor $\cons_{A}$ in $\Fam(P)$, one
rule with $\mathrm{ar}(A)$ loud premises at slots $0..\mathrm{ar}(A)-1$ and
conclusion \code{Encode\{deliver: Slot(ar(A)), node: Atom(A, payloads)\}};
and under \code{inst} the rule of Definition~\ref{def:inst}, one loud premise
at slot 0 and conclusion \code{Encode\{deliver: Slot(1), node:
Instantiate\{ctx: Slot(2), with: Payload(0)\}\}}. Rule classes, join plans,
the grading and $w$ \MUST{} be derived from the emitted table.
\end{requirement}

\begin{requirement}[Only a forwarder reads a store]\label{req:quietjoin}
As v0.1, and now a property of the generator: no generated rule other than
$\mathsf{quote}$ \MAY{} have a quiet premise. $\cons_{\qq}$ \emph{builds} a
store name and reads only messages. A test \MUST{} assert, for every generated
consumer shape other than $\fw$, that it has no redex against a $\qq$ row.
\end{requirement}

\begin{requirement}[Classes and plans are checked]\label{req:plans}
A test \MUST{} check, for the generated table of every corpus artefact: exactly
one class-(c) rule; one class-(b) rule per member of $\Fam(P)$, plus
$\mathsf{make}_{\mid}$ and, under \code{inst}, $\inst$; seven class-(a)
rules; and a join plan per rule whose arity equals its premise count. This
replaces v0.1's hand-stated count, which the family made obsolete.
\end{requirement}

\subsection{The grading, generated}

\begin{requirement}[Grading]\label{req:grading}
The weight $\Wt$ of \cite[Lem.~9.2]{draft3} --- $\bl,\br,\sy$ at two, $\mm$ at
zero, every other non-message shape at one --- \MUST{} be derived from the
table as in v0.1: one plus the heaviest non-message atom any of a shape's
rules produces from it. For the family and for $\inst$ this gives one. The
debug build \MUST{} assert that every non-$\mathsf{eval}$ firing decreases
$\Wt$, and that without $\dd$ or $\ev$ firings a step does not increase
messages plus stores (\cite[Thm.~6.1(iii)]{matnote}).
\end{requirement}

%======================================================================
\section{Sparse algebra and redex finding}\label{sec:finders}

The sparse layer (\code{Semiring}, \code{Csr}, Gustavson \code{spgemm},
counting-sort \code{bucket}) is unchanged from v0.1, as are the three finders
(\code{Naive}, \code{Product}, \code{Bucketed}), the canonical key, the
distinctness and symmetry masks, and the requirement that all three return
identical sorted redex sets. Two things change.

\begin{requirement}[Joins are $n$-way]\label{req:nway}
A join plan has up to $w-1$ premises. The \code{Product} finder \MUST{}
compute one product per premise slot and an $n$-fold row-wise expansion with
the masks; the \code{Bucketed} finder \MUST{} look up one bucket per slot.
Neither \MAY{} special-case arity. The redex type is
\code{(rule: u16, consumer: u32, premises: [u32; W-1])}.
\end{requirement}

\begin{requirement}[Validation covers the generated signature]\label{req:validate}
Finder agreement (Obligation~\ref{obl:finders}) \MUST{} be run over every
shape the generator can emit for the corpus, two levels of the family, and
$\inst$ --- the reference's version~1.2 validation does this over
twenty-nine shapes (\cite[Find.~10.1]{matnote}).
\end{requirement}

\begin{hazard}[Contention at static channels]\label{haz:statics}
v0.1 flagged the quadratic join on a contended channel
(\cite[Rem.~7.2]{matnote}). Phase two makes static channels contended by
design: $n$ instances put $n$ messages at each allocation channel under
literal erection, and at $r$ under \code{inst}. Under \code{inst} the
contended rule has one premise and the redex count at $r$ is $n^{2}$, whose
maximal independent set is a perfect matching, as in E2; under literal
erection the multi-premise constructors at static channels enumerate
$n^{m}$ candidates for an $m$-input member. The enumerated-to-fired ratio
\MUST{} be reported per static channel.
\end{hazard}

%======================================================================
\section{Conflict, selection, and resolvers}\label{sec:select}

The conflict hypergraph, the priority function, the greedy maximal independent
set and the proposition that greedy selection equals parallel rounds under
shared priorities are unchanged from v0.1 (\code{mix64} of SplitMix64 over
seed, step and canonical key; \cite{bfs}). So is the requirement that every
backend compute exactly the greedy set.

\subsection{Resolvers}\label{sec:resolvers}

\begin{lstlisting}
pub enum Resolver {
    MaximalProgress,     // the greedy maximal step; the default
    SingleByPriority,    // the first redex in priority order; linearises a step
    SingleUniform,       // one redex, uniformly over the (symmetry-masked) redex set
    CrossInstance,       // maximal progress, preferring redexes that join provenances
}
\end{lstlisting}

\begin{requirement}[The cross-instance resolver]\label{req:cross}
\code{CrossInstance} \MUST{} compute the greedy maximal independent set under
the priority $(\,\neg x(r),\ h(r),\ \mathrm{key}(r)\,)$, where $x(r)$ holds
iff the premises of $r$ carry two distinct non-\code{STATIC} provenance tags.
It is the matrix machine's counterpart of the adversarial profile
\cite[Obl.~11.2]{combspec} requires of \Camp, and is stronger: it prefers
every cross at every static channel in the same step. It \MUST{} be used in
the phase-two tests of Section~\ref{sec:test} alongside
\code{MaximalProgress}, and a term that passes \Comb's conformance check
\cite[Req.~5.14]{combspec} yet yields a \code{MIXED} row under it \MUST{} be
reported as a counterexample to the conformance check, not patched.
\end{requirement}

\begin{requirement}[Races are visible]\label{req:races}
As v0.1, with \cite[Req.~9.8]{combspec}'s third race added: every selected
redex at a static channel whose excluded rivals carry a different provenance
\MUST{} be a trace event at level \code{races}, naming both addresses by
content hash.
\end{requirement}

%======================================================================
\section{Firing}\label{sec:fire}

\begin{requirement}[Order of effects]\label{req:insertorder}
As v0.1: sort the step by canonical key; compute conclusions in that order,
class (b) interning new names in that order --- including the names an
$\inst$ substitution creates, in the order of a fixed traversal of $C$ ---
and class (c) reading the arena; remove consumed rows by stable compaction;
append produced rows in the order computed. Parallel backends \MAY{} compute
concurrently but \MUST{} assign ids as if sequentially. Budget checks precede
any mutation.
\end{requirement}

\begin{requirement}[Soundness is checked, not assumed]\label{req:soundcheck}
As v0.1: the debug build re-finds redexes after each firing of a step in
canonical order and checks the rest remain enabled.
\end{requirement}

%======================================================================
\section{Addresses and freshness}\label{sec:addr}

\begin{requirement}[One allocator for both machines]\label{req:oneallocator}
The address allocator of \cite[Req.~9.5, Req.~5.23]{combspec} \MUST{} be
factored out of \code{f1r3comb-run} into \code{f1r3comb-term} (module
\code{addr}) and used unchanged by \Mat, so that a deploy on either machine
receives the same root address for the same program and allocator state. Two
machines that disagree on addresses cannot be compared trace for trace.
\end{requirement}

\begin{requirement}[Deploy checks]\label{req:deploychecks}
Deploy \MUST{} perform \Comb's three freshness checks, with its diagnostics:
the root address exceeds, by encoding size, every name of the image
(\code{comb-address-too-small}, \cite[Req.~5.19]{combspec}); so does every
static channel (\code{comb-static-not-fresh}, \cite[Req.~5.22]{combspec});
and the root address is prefix-incomparable with every address in use
(\code{comb-address-overlap}, \cite[Req.~5.23]{combspec}). Size comparison
uses the \code{size} column; prefix comparison walks the \code{root} and
$\Lft/\Rgt$ spine and is linear in depth.
\end{requirement}

%======================================================================
\section{The operator layer, scoped}\label{sec:op}

The operator layer of v0.1 --- species, markings, successors with
mass-action coefficients, exploration, the materialised $T$, triangularity in
the marking order graded by $\Wt$, closure by repeated squaring, and
$T^{\top}\chi$ --- is unchanged in content. Its scope is restated.

\begin{requirement}[Scope, stated in the tool]\label{req:opscope}
Every phase-two image contains $\ev$ (each gate stores its body behind one
\cite[Lem.~6.2]{draft3}, and each erection ends in one), so
\code{explore} \MUST{} refuse every compiled program with an input with
\code{comb-op-destructor} unless given a depth bound, and the documentation
\MUST{} say that triangularity and closure
(\cite[Thm.~6.1(ii), Cor.~6.2]{matnote}) apply to hand-written terms and to
phase-one images of input-free programs only (\cite[\S14.2]{matnote}). The
\code{successors} primitive is unrestricted and is what the logic's modality
needs at a frontier.
\end{requirement}

\begin{remark}[\Comb's open question on the matrix encoding]
\cite[\S15, item~5]{combspec} observes that an upper-triangular encoding wants
a bounded name universe and that phase two grows it with every unfolding. The
answer, from this side: the machine never wanted a bounded universe --- it
multiplies the term-sized incidence, not the state-indexed operator
(\cite[\S7]{matnote}), and its per-unfolding name cost is constant --- and the
closure, which does want one, is unavailable to compiled programs with inputs
for the independent reason that they contain $\ev$.
\end{remark}

%======================================================================
\section{Instrumentation and the width gate}\label{sec:instr}

\begin{requirement}[Reports]\label{req:reports}
As v0.1 (steps, rewrites, width and enumeration histograms, names, arena,
inert rows, per-rule counts), plus: steps and rewrites attributed to
\emph{erection} --- $\inst$ and constructor firings whose conclusion is
consumed by $\ev$ in the next step, and those $\ev$ firings --- separately
from the rest; instances started; address depth high-water mark; new names per
instance; and the count of \code{MIXED} rows.
\end{requirement}

\begin{requirement}[The width measurement]\label{req:widthcmd}
\code{f1r3x comb-width} \MUST{} report, per program and per erection, the
width distribution with erection separated from computation, as
\cite[Obl.~14.6]{matnote} asks. Under \code{inst}, erection is two steps per
instance and parallel across instances, so it should add width and not depth;
the report is where that is confirmed or refuted. This is the instrument for
\cite[Obl.~13.7]{matnote}, which \cite[Obl.~14.5]{matnote} now gates.
\end{requirement}

\begin{requirement}[Subcommands]\label{req:fx}
\Fx{} gains \code{comb-mat}, \code{comb-run --machine=matrix} (with
\code{-{}-resolver}, \code{-{}-erection}, \code{-{}-seed}, \code{-{}-cores},
\code{-{}-trace}), \code{comb-width} and \code{comb-op}, alongside \Comb's
\code{comb}, \code{comb-ir}, \code{comb-hash}, \code{comb-run} and
\code{comb-lattice}.
\end{requirement}

%======================================================================
\section{Architecture}\label{sec:arch}

\begin{center}
\begin{tabularx}{\linewidth}{@{}l>{\raggedright\arraybackslash}Xl@{}}
\toprule
Crate & What it is & Depends on\\
\midrule
\code{f1r3comb-term} & (\Comb's) plus the rule generator
(Decision~\ref{dec:ruletable}) and the address allocator
(Requirement~\ref{req:oneallocator}) & \code{k1ndl1ng-norm}\\
\code{f1r3comb-mat} & Intern table with size, root and depth columns; rows of
width $w$; shape tables with provenance; import, export, \code{.cmat}; sparse
layer & \code{f1r3comb-term}\\
\code{f1r3comb-par} & Join plans, finders, conflict, priority, resolvers,
firing (including $\inst$), deploy checks, backends, reports &
\code{f1r3comb-mat}\\
\code{f1r3comb-op} & Species, markings, successors, exploration, $T$,
closure, \code{pre} & \code{f1r3comb-par}\\
\code{f1r3comb-check} (shared) & Python parity, trace certification,
phase-two tests, differential tests against \code{f1r3comb-run} and \Bell &
all\\
\bottomrule
\end{tabularx}
\end{center}

\begin{requirement}[Independence]\label{req:indep}
As v0.1, and in particular: none of the three crates \MAY{} depend on
\code{f1r3comb-ir}, \code{f1r3comb-lower} or \code{f1r3comb-run}.
\end{requirement}

\subsection{Backends and the device seam}\label{sec:device}

The \code{Backend} trait, the sequential and threaded host backends, and the
rule that device backends live outside the workspace are unchanged from v0.1.
The kernel contract is restated for the generated table.

\begin{requirement}[Device kernels]\label{req:device}
A device backend \MUST{} implement, against the \code{.cmat} layout: bucketing
as a sort by key; one segmented $n$-way expansion per join plan with the
masks; \code{mix64} priorities; the rounds of the greedy-equals-rounds
proposition; stable compaction and a prefix-sum append; class (b) as a batched
insert-or-get; class (c) as a gather from the arena; and, under \code{inst},
substitution as a gather over the context's interned form with a batched
insert-or-get of the results. It \MUST{} target $w = 4$ and the \code{inst}
erection (Decision~\ref{dec:twoerect}); a device backend for literal erection
\MUSTNOT{} be built, because under concurrency it computes wrong answers
(Hazard~\ref{haz:gate}). It is accepted only on identical step traces to
\code{HostSequential} on the conformance corpus, and no work on it \MAY{}
begin before Obligations~\ref{obl:inst} and~\ref{obl:width} are discharged.
\end{requirement}

%======================================================================
\section{Correctness obligations and testing}\label{sec:test}

\begin{obligation}[Carried from v0.1]\label{obl:carried}
Round trip; finder agreement over the generated signature
(Requirement~\ref{req:validate}); Python parity on E1--E4, exact against
\code{results.txt}; trace certification by an independent single-redex
rewriter written from \cite[Def.~4.4]{combspec}; operator correctness on
destructor-free terms; determinism across backends, core counts and
WebAssembly; the invariants of Requirement~\ref{req:grading}; device
conformance corpus.
\end{obligation}

\begin{obligation}[Finder agreement]\label{obl:finders}
Requirement~\ref{req:validate}: on at least $10^{4}$ random states over every
generated shape, the three finders agree exactly.
\end{obligation}

\begin{obligation}[The $W$ regression, on this machine]\label{obl:W}
\cite[Obl.~11.1]{combspec}'s term $W$ and its three-instance variant, compiled
by \Comb, run under \Mat{} with \code{MaximalProgress} and
\code{CrossInstance}, for 100 seeds each, under each erection. Under
\code{inst} the observation at $\enc{\quo{X}}$ \MUST{} match \Bell's on every
seed and no row \MUST{} be \code{MIXED}. Under literal erection the test
\MUST{} record the fraction of seeds producing the crossed residue; a
fraction of zero under \code{CrossInstance} would contradict
\cite[Find.~14.1]{matnote} and \MUST{} be investigated.
\end{obligation}

\begin{obligation}[The phase-two experiments, reproduced]\label{obl:phase2}
\code{f1r3comb-check} \MUST{} reproduce \code{phase2-results.txt} in kind:
the mixing counts of \cite[Find.~14.1]{matnote} are seed-dependent and are
compared as rates within a tolerance; the $D_{x}$ laws of
Requirement~\ref{req:dxlaw} are exact.
\end{obligation}

\begin{obligation}[Addresses]\label{obl:addr}
\cite[Obl.~11.5, Obl.~11.10]{combspec} on this machine: the three deploy
checks, including the adversarial source term $\quo{(\out{x}{\nil})}$ whose
encoding is a leaf of $\enc{\quo{x}}$; and address independence, by running
the corpus at two unrelated root addresses and comparing observations.
\end{obligation}

\begin{obligation}[Differential testing]\label{obl:diff}
For programs whose observations are resolver-independent, \Mat, \code{f1r3comb-run}
and \Bell{} \MUST{} agree on the observation of \cite[Obl.~11.7]{combspec}.
Under \code{inst} this requires \Comb{} to emit the \code{inst} erection
(Section~\ref{sec:paper}, item~2); until it does, \Mat{} compiles
\code{inst} images from \Comb's phase-one output with its own test-only
emitter in \code{f1r3comb-check}, which \MUSTNOT{} be linked into any
product crate.
\end{obligation}

\begin{obligation}[The width of real programs]\label{obl:width}
\cite[Obl.~13.7]{matnote}, under \code{inst}, with erection separated. Gates
the device seam.
\end{obligation}

\begin{obligation}[Phase two's primitive]\label{obl:inst}
\cite[Obl.~14.5]{matnote}: the paper decides phase two's primitive. Gates
correctness of compiled programs under concurrency, hence the device seam.
\end{obligation}

%======================================================================
\section{Errors and diagnostics}\label{sec:errors}

Shared with \Comb{} \cite[\S12]{combspec}: \code{comb-feature-mismatch},
\code{comb-presentation-mismatch}, \code{comb-name-budget},
\code{comb-address-too-small}, \code{comb-static-not-fresh},
\code{comb-address-overlap}, \code{comb-decode}. Specific to \Mat:

\begin{center}
\begin{tabularx}{\linewidth}{@{}ll>{\raggedright\arraybackslash}X@{}}
\toprule
Code & Where & Meaning\\
\midrule
\code{comb-mat-presentation} & any & Built or invoked under
\code{presentation-b}\\
\code{comb-mat-width} & import & The signature needs rows wider than the
instantiated $w$\\
\code{comb-mat-hole} & import, step & \code{HOLE} outside the context
argument of an $\inst$ atom, or in a conclusion\\
\code{comb-mat-hash-collision} & boundary & Structural key and content hash
disagree\\
\code{comb-mat-overflow} & sparse, op & Checked \code{Nat} arithmetic\\
\code{comb-mat-decode} & import & Malformed \code{.cmat}\\
\code{comb-op-destructor} & op & Exploration of a term that can reach $\ev$,
without a depth bound; every compiled program with an input\\
\code{comb-op-budget} & op & Exploration budget exceeded\\
\code{comb-mat-internal} & any & An invariant assertion failed; names the
requirement\\
\bottomrule
\end{tabularx}
\end{center}

%======================================================================
\section{Performance, portability, determinism}\label{sec:perf}

\begin{requirement}[Complexity]\label{req:complexity}
As v0.1; in addition a step \MUSTNOT{} depend on address depth, which
structural interning secures, and an $\inst$ firing \MUST{} cost $O(|C|_{\#})$
where $|C|_{\#}$ is the number of distinct names in the context.
\end{requirement}

\begin{requirement}[Targets]\label{req:targets}
v0.1's budget figures, plus: $10^{4}$ unfoldings of $D_{x}$ under \code{inst}
with no rise in per-unfolding step time, which is \cite[Req.~13.2]{combspec}'s
test of prefix sharing met by construction; and 64 concurrent instances of a
$10^{3}$-atom unit erected in two steps.
\end{requirement}

\begin{requirement}[Portability and determinism]\label{req:port}
As v0.1 and \cite[Req.~13.3]{combspec}: \code{wasm32-unknown-unknown},
\code{forbid(unsafe\_code)}, no external dependency, no \code{RandomState};
bit-identical traces, reports and exports across runs, hosts, core counts and
WebAssembly, given the same term, root address, seed, resolver, erection and
configuration.
\end{requirement}

%======================================================================
\section{What the implementation asks of the papers}\label{sec:paper}

\paragraph{Discharged since v0.1.} The atom-count lemma is weighted and the
length bound restated \cite[Lem.~9.2, Prop.~9.9]{draft3}; the matrix note's
nilpotency index, closure and message-plus-store invariant are corrected
(\cite[Thm.~6.1, Cor.~6.2]{matnote}); the reference counts a symmetric match
once; the rules of the shortcut constructors are now members of the family
\cite[Def.~3.8]{draft3}; the rule table is generated
\cite[Req.~4.5]{combspec}; the $\cons_{\qq}$ reconciliation v0.1 asked for is
\cite[Rem.~3.9]{draft3}, and is trivial in the row representation.

\paragraph{Open.}
\begin{enumerate}[leftmargin=1.8em]

\item \textbf{Draft~3, Definition~6.10 and Remark~6.12.} The discipline is
stated for constructors, but a router holding a leaf at a static subject
joins instances just as well ($\fw(k,\Lft\sigma) \mid \mm(k,\Lft\tau)$), so
the proof of Lemma~6.11 needs the discipline for every redex
(\cite[Def.~14.2]{matnote}). In that form it cannot be met by fixed-arity
constructors for any atom mentioning two allocated names
(\cite[Prop.~14.3]{matnote}), so Remark~6.12 describes no construction, and
Remark~3.14's claim that the context-indexed constructor has ``the same
reach'' as the family is true of reachability but false under the
discipline. Proposed amendment: promote the context-indexed constructor, in
its single-input multi-hole form, to a primitive, and state phase two with
it. \textsf{GATE}

\item \textbf{\Comb{} v0.4, Requirement~5.16.} The two-stage erection's
stage~(ii) is unreachable from stage~(i) by the same proposition, so the
requirement mandates a construction that does not exist. Proposed amendment:
replace it with the $\inst$ erection of Definition~\ref{def:inst}, which
removes Requirements~5.9, 5.12 and~5.16 and the constructors of constructors,
and keeps 5.11, 5.19--5.23 unchanged.

\item \textbf{\Comb{} v0.4, the tag table of Requirement~4.11.} The
constructor flag covers constructors of base shapes, but under literal
erection $\Fam(P)$ contains $\cons_{\cons_{A}}$ whenever the program has an
(o5) occurrence, and those have no tag. Moot under item~2; otherwise the
layout needs a level field.

\item \textbf{\Comb{} v0.4, tags for the context constructor.} Under item~2,
$\inst$ needs a shape tag and the hole needs a reserved name tag outside the
source block (Requirement~\ref{req:hole}).

\item \textbf{\Comb{} v0.4, Requirement~9.5.} It asks for ``a root address of
a shape outside that program's family'', which is v0.3's shape-based
freshness; Requirement~5.19 and draft~3's Lemma~6.14 are by size. Proposed
amendment: restate 9.5 by size.

\item \textbf{\Comb{} v0.4, citations to the logic note.} The weighted count
is cited as \cite{logic}'s ``Def.~6.6, Prop.~9.2''; it is Definition~4.2 and
Proposition~6.17 there, and \cite[Lem.~9.2]{draft3} in the companion.

\item \textbf{\Comb{} v0.4, Requirements~9.7 and~13.1.} Prefix sharing makes
\emph{interning} $O(1)$ per level, but the content hash of the full encoding
remains $O(\text{depth})$, and \Camp{} keys channels by it. Proposed: a
Merkle form of the canonical hash (tag, then child hashes) for names, which
makes a new address's hash $O(1)$ and keeps equality of hashes equivalent to
$\nequiv$; it is a change to \cite{k1ndl1ng}'s hashing and wants deciding
there.

\item \textbf{\Comb{} v0.4, Requirement~9.3} (decode by projection),
unchanged from v0.1: scope it to builds where \code{presentation-b} is
reachable (Decision~\ref{dec:decode}).

\item \textbf{\Comb{} v0.4, the address allocator.} Move it to
\code{f1r3comb-term} so both machines share it
(Requirement~\ref{req:oneallocator}).

\end{enumerate}

%======================================================================
\section{Open questions}

\begin{enumerate}[leftmargin=1.8em]
\item \textbf{The primitive} (Obligation~\ref{obl:inst}): the single-input
context constructor, or another construction meeting the discipline.
\item \textbf{Width under instantiation.} Whether erection, at two steps per
instance, adds width without depth on real programs.
\item \textbf{Resolver fidelity and fairness} (\cite[Obl.~13.1,
Obl.~13.6]{matnote}); now also: which resolver a correct phase-two image is
insensitive to, which the permutation lemma says is all of them.
\item \textbf{Contention at statics} (Hazard~\ref{haz:statics}): under
\code{inst} the quadratic join sits at $r$, one per unit; whether bucketed
matching at a one-premise channel should bypass enumeration.
\item \textbf{Address compaction} \cite[\S15, item~1]{combspec}: whether a
retired subtree may be reused, which the intern table would make cheap and
the permutation lemma would have to license.
\item \textbf{Garbage}, \textbf{the history unfolding}, \textbf{weights}: as
v0.1.
\end{enumerate}

%======================================================================
\section{Delivery}

\begin{center}
\begin{tabularx}{\linewidth}{@{}ll>{\raggedright\arraybackslash}X@{}}
\toprule
M & Deliverable & Done when\\
\midrule
X0 & \code{f1r3comb-mat} & Intern table with size, root and depth columns;
rows of width $w$; provenance; import, export, \code{.cmat}; round trip.
Needs \Comb{} M0 and the rule generator in \code{f1r3comb-term}\\
X1 & Finders & $n$-way plans from the generated table; three finders; masks;
agreement over the generated signature\\
X2 & \code{f1r3comb-par}, sequential & Resolvers including
\code{CrossInstance}; firing including $\inst$; deploy checks with the shared
allocator; Python parity; certification; invariants; the $D_{x}$ laws\\
X3 & Phase-two evidence & The $W$ regression under both erections and both
resolvers; the mixing rates of \cite[Find.~14.1]{matnote}; reported to the
paper as evidence for Obligation~\ref{obl:inst}\\
X4 & The width gate & Under \code{inst}, with erection separated; needs
\Comb{} M2 or the test-only emitter. \textbf{Gate} on X7\\
X5 & Threaded host & Determinism across cores and \code{wasm32};
differential tests\\
X6 & \code{f1r3comb-op} & As v0.1, scoped per
Requirement~\ref{req:opscope}\\
X7 & Device conformance & Corpus under \code{inst}; external backend only
after Obligations~\ref{obl:inst} and~\ref{obl:width}\\
\bottomrule
\end{tabularx}
\end{center}

\noindent
X3 is new and comes early on purpose. The machine is the cheapest instrument
there is for the question phase two now hinges on, because maximal progress
and \code{CrossInstance} produce, as a matter of course, the interleavings a
store-based machine must be coaxed into.

%======================================================================
\section{Changes from v0.1}\label{sec:v1delta}

\begin{center}
\begin{tabularx}{\linewidth}{@{}l>{\raggedright\arraybackslash}X@{}}
\toprule
v0.1 & v0.2\\
\midrule
Thirteen shapes, fifteen with shortcuts; fixed width four & Generated
signature; width $w$ per artefact, four under \code{inst}\\
One rule table, shared & Generated from \Comb's shape table\\
Shortcut constructor rules fixed here & Members of the family\\
No addresses & \Comb's allocator, shared; size, root and depth columns;
deploy checks; provenance\\
Three resolvers & Four; \code{CrossInstance} added\\
--- & Section~\ref{sec:phase2}: phase two as the machine meets it; the
$\inst$ erection behind a feature; Hazard~\ref{haz:gate}\\
Operator layer for the $\ev$-free sublattice & The same, with its scope
stated: no compiled program with an input\\
Three errata against the notes & Discharged by their revisions\\
Eleven asks of the papers & Nine, all but one new, the first a gate\\
Delivery X0--X6 & X0--X7, with the phase-two evidence at X3\\
\bottomrule
\end{tabularx}
\end{center}

%======================================================================
\begin{thebibliography}{99}

\bibitem{matnote} L.\,G. Meredith.
\emph{Rho combinators as sparse linear algebra.} Version~1.2,
\code{publications/rho-combinators/rhocomb-matrix}, with \code{rhocombmat.py},
\code{experiments.py} and \code{phase2.py}, September 2026.

\bibitem{combspec} F1R3FLY.io.
\emph{F1R3Comb: a compiler from kernel rho to the name-free rho combinators,
and a machine that runs the result.} Specification v0.4,
\code{publications/rho-combinators/f1r3comb-spec-v04}, 23 September 2026.

\bibitem{draft3} L.\,G. Meredith and M. Stay.
\emph{Name-free combinators for the rho calculus: a repair, a linear encoding,
and the one operation that was missing.} Draft~3,
\code{publications/rho-combinators}, September 2026.

\bibitem{logic} F1R3FLY.io.
\emph{A spatial-behavioural logic for the rho combinators.} Draft~2,
revision~2, \code{publications/rho-combinators/rhocomb-logic}, 23 September
2026.

\bibitem{k1ndl1ng} F1R3FLY.io.
\emph{K1ndl1ng: a parser and normaliser for kernel rholang.}
\code{publications/campf1r3/k1ndl1ng-spec}, 2026.

\bibitem{bfs} G.\,E. Blelloch, J.\,T. Fineman and J. Shun.
Greedy sequential maximal independent set and matching are parallel on
average. \emph{SPAA}, 2012.

\bibitem{gustavson} F.\,G. Gustavson.
Two fast algorithms for sparse matrices: multiplication and permuted
transposition. \emph{ACM Transactions on Mathematical Software}, 1978.

\bibitem{splitmix} G.\,L. Steele, D. Lea and C.\,H. Flood.
Fast splittable pseudorandom number generators. \emph{OOPSLA}, 2014.

\end{thebibliography}

\end{document}
